% unidades/14-te-form-2.tex
\chapter{🔗 て形 (2) — Permisos y Prohibiciones}
\unidadheader{14}{て形 (2)}{La Forma て}{Hablar de lo que se puede y no se puede hacer}

\begin{tcolorbox}[vocabbox, title={📌 Vocabulario Nuevo}]
\begin{tabularx}{\linewidth}{llX}
\toprule
\textbf{Japonés} & \textbf{Español} & \textbf{Tipo} \\
\midrule
\vocabitem{写真を撮る}{しゃしんをとる}{tomar una foto}{v. gr.1}
\vocabitem{タバコを吸う}{たばこをすう}{fumar}{v. gr.1}
\vocabitem{車を止める}{くるまをとめる}{estacionar el auto}{v. gr.2}
\vocabitem{遊ぶ}{あそぶ}{jugar / divertirse}{v. gr.1}
\vocabitem{急ぐ}{いそぐ}{darse prisa}{v. gr.1}
\vocabitem{忘れる}{わすれる}{olvidar}{v. gr.2}
\vocabitem{困る}{こまる}{estar en problemas}{v. gr.1}
\vocabitem{休む}{やすむ}{descansar / faltar}{v. gr.1}
\bottomrule
\end{tabularx}
\end{tcolorbox}

\subsection*{🗣️ Frases Clave}

\frase{ここで\ruby{写真}{しゃしん}を\ruby{撮}{と}ってもいいですか？}{¿Puedo tomar una foto aquí?}
\frase{すみません、ここはダメです。}{Lo siento, aquí no se puede.}
\frase{ここでタバコを\ruby{吸}{す}ってはいけません。}{No se debe fumar aquí.}
\frase{\ruby{一緒}{いっしょ}に\ruby{遊}{あそ}びましょう！}{¡Vamos a jugar juntos!}

\subsection*{⚙️ Gramática}

\patron{Pedir permiso: 〜てもいいですか}
  {V-て + もいいですか}
  {\ej{\ruby{入}{はい}ってもいいですか？}{¿Puedo entrar?}
  \ej{これを\ruby{使}{つか}ってもいいですか？}{¿Puedo usar esto?}}

\patron{Prohibición: 〜てはいけません}
  {V-て + はいけません}
  {\ej{ここに\ruby{車}{くるま}を\ruby{止}{と}めてはいけません。}{No se debe estacionar el auto aquí.}
  \ej{お\ruby{酒}{さけ}を\ruby{飲}{の}んではいけません。}{No se debe beber alcohol.}}

\begin{tcolorbox}[ejerciciobox, title={📝 Ejercicios Rápidos}]
\begin{enumerate}
  \item Pide permiso para sentarte.
  \item Di que no se debe entrar ahí.
  \item Traduce: \ruby{教}{きょう}\ruby{室}{しつ}で\ruby{寝}{ね}てはいけません。
\end{enumerate}
\vspace{0.4em}
\textbf{Respuestas:}
1) \ruby{座}{すわ}ってもいいですか？\quad
2) そこに\ruby{入}{はい}ってはいけません。\quad
3) No se debe dormir en el salón de clases.
\end{tcolorbox}

\begin{tcolorbox}[kanjibox, title={🀄 Kanjis de la Unidad}]
\begin{tabularx}{\linewidth}{cllXX}
\toprule
\textbf{Kanji} & \textbf{On} & \textbf{Kun} & \textbf{Significado} & \textbf{Vocab} \\
\midrule
\kanjirow{使}{し}{つか}{usar}{\ruby{使}{つか}う / \ruby{大使館}{たいしかん}}
\kanjirow{置}{ち}{お}{poner / colocar}{\ruby{置}{お}く / \ruby{設置}{せっち}}
\kanjirow{急}{きゅう}{いそ}{apresurarse}{\ruby{急}{いそ}ぐ / \ruby{急行}{きゅうこう}}
\kanjirow{休}{きゅう}{やす}{descansar}{\ruby{休}{やす}む / \ruby{休日}{きゅうじつ}}
\kanjirow{始}{し}{はじ}{comenzar}{\ruby{始}{はじ}まる / \ruby{開始}{かいし}}
\bottomrule
\end{tabularx}
\end{tcolorbox}
